\section{Body-Space Direct Solving Utilizing the Woodbury Matrix Identity}
\label{sec:body-space}

\subsection{Constraint Space and Body Space}
\label{sec:body-space-idea}

The system \eqref{eq:ostrich-system} is one of two ways to write the
linear system of a Newton iteration. Both come from a block system in two
unknowns, the velocity update $\Delta\mathbf{u}$ of the bodies and the
constraint force update $\Delta\boldsymbol{\lambda}$:
\begin{equation}
  \begin{bmatrix}
    \tilde{\mathbf{M}} & -\mathbf{J}^\top \\
    \mathbf{J} & \mathbf{C}
  \end{bmatrix}
  \begin{bmatrix} \Delta\mathbf{u} \\ \Delta\boldsymbol{\lambda} \end{bmatrix}
  = -\begin{bmatrix} \mathbf{r}_u \\ \mathbf{r}_\lambda \end{bmatrix} ,
  \label{eq:kkt}
\end{equation}
where $\mathbf{r}_u = \mathbf{r}_{\text{dyn}}/h$ and
$\mathbf{r}_\lambda = \mathbf{r}_{\text{con}}/h$, so that eliminating
$\Delta\mathbf{u}$ gives exactly \eqref{eq:ostrich-system}. Ostrich uses
maximal coordinates: every rigid body has $6$ velocity degrees of freedom,
so $\Delta\mathbf{u} \in \mathbb{R}^{6N_b}$ for $N_b$ bodies. The
generalized mass matrix $\tilde{\mathbf{M}}$ is block diagonal, with one
$6 \times 6$ block (mass and inertia) per body. The Jacobian
$\mathbf{J} \in \mathbb{R}^{n \times 6N_b}$ has one row per constraint, and
each row acts on at most two bodies. The compliance $\mathbf{C} =
\operatorname{diag}(c_1, \dots, c_n)$ is diagonal.

The block system can be reduced to either unknown:
\begin{itemize}
  \item \textbf{Constraint space} (ostrich today). The first block row gives
    $\Delta\mathbf{u} = \tilde{\mathbf{M}}^{-1}(\mathbf{J}^\top\Delta\boldsymbol{\lambda} - \mathbf{r}_u)$,
    which is cheap because $\tilde{\mathbf{M}}$ is inverted block by block.
    Substituting it into the second row gives \eqref{eq:ostrich-system}. The
    unknowns are the constraint forces $\Delta\boldsymbol{\lambda}$, and the
    system has dimension $n$.
  \item \textbf{Body space}. The second block row gives
    $\Delta\boldsymbol{\lambda} = -\mathbf{C}^{-1}(\mathbf{r}_\lambda + \mathbf{J}\Delta\mathbf{u})$,
    which is cheap because $\mathbf{C}$ is diagonal. Substituting it into the
    first row gives
    \begin{equation}
      \underbrace{\left(\tilde{\mathbf{M}} + \mathbf{J}^\top\mathbf{C}^{-1}\mathbf{J}\right)}_{\mathbf{K}}
      \Delta\mathbf{u}
      = -\mathbf{r}_u - \mathbf{J}^\top\mathbf{C}^{-1}\mathbf{r}_\lambda .
      \label{eq:body-system}
    \end{equation}
    The unknowns are the body velocities $\Delta\mathbf{u}$, and the system
    has dimension $6N_b$. The forces then follow from the formula for
    $\Delta\boldsymbol{\lambda}$ above.
\end{itemize}
Both reductions are Schur complements of \eqref{eq:kkt} and give the same
solution. What changes is the dimension, from $n$ to $6N_b$. For Helhest,
$N_b = 4$, so the body-space system is always $24 \times 24$, while
$n = 31$ to $67$ in the obstacle scene and $50$ to $170$ in the terrain
scene (Section~\ref{sec:choosing-z}). The body-space dimension depends only
on the number of bodies. It does not grow with the number of active
contacts, each of which adds three rows (one normal, two friction) to the
constraint-space system.

\paragraph{Prerequisite: $c_i > 0$.}
The body-space form needs $\mathbf{C}^{-1}$, i.e.\ $c_i > 0$ for every
active row. A perfectly rigid constraint ($c_i = 0$) has no body-space
form. In ostrich every row is compliant. In the recorded Helhest systems the
compliance is $0.011$ for joint rows, $0.22$ for control rows, $0.11$ to
$0.95$ for contact normals and $10^{-6}$ to $3 \cdot 10^{6}$ for friction
rows, and the system operator adds a regularization of $10^{-6}$ to every
diagonal entry. Very small $c_i$ are still a numerical concern: they put
weights $1/c_i$ of order $10^{6}$ into $\mathbf{K}$, which limits the
accuracy of a single-precision solve.

\paragraph{A direct solver in body space.}
$\mathbf{K}$ is SPD, because $\tilde{\mathbf{M}}$ is SPD and
$\mathbf{J}^\top\mathbf{C}^{-1}\mathbf{J}$ is positive semidefinite. A
Cholesky factorization \eqref{eq:cholesky} therefore applies, without
pivoting. Its cost is about $\tfrac{1}{3}(6N_b)^3$ operations
(Table~\ref{tab:direct-cost}): about $4.6 \cdot 10^3$ for Helhest, against
about $1.6 \cdot 10^6$ for a dense Cholesky factorization of $\mathbf{A}$
at $n = 170$. Assembling $\mathbf{K}$ is cheap too. Each row adds
$c_i^{-1}\mathbf{j}_i^\top\mathbf{j}_i$, where $\mathbf{j}_i$ is the $i$-th
row of $\mathbf{J}$, and this term touches at most two body blocks, so the
assembly costs $\mathcal{O}(n)$. The small, fixed size is what makes an
exact solve affordable in every Newton iteration. The soft modes that slow
down PCR are no obstacle here. In the directions of
$\operatorname{null}(\mathbf{J}^\top)$, where
$\mathbf{A}\mathbf{x} = \mathbf{C}\mathbf{x}$
(Section~\ref{sec:choosing-z}), the body-space form handles the forces
exactly through the division by $c_i$. In our offline measurements, this
$24 \times 24$ solve is $3.9$ to $5.6$ times faster than the best PCR
variant for a single simulation, and $1.7$ to $2.0$ times faster for $1024$
parallel simulations.

The argument relies on $6N_b$ being small. With many bodies, a dense
factorization of $\mathbf{K}$ grows like $(6N_b)^3$. $\mathbf{K}$ is sparse
(the block of bodies $i$ and $j$ is nonzero only if a constraint row couples
them), but for large connected piles of bodies PCR is expected to win again.

\subsection{Relation to the Woodbury Identity}
\label{sec:woodbury}

The ostrich system matrix
$\mathbf{A} = \mathbf{J}\tilde{\mathbf{M}}^{-1}\mathbf{J}^\top + \mathbf{C}$
is a diagonal matrix $\mathbf{C}$ plus the term
$\mathbf{J}\tilde{\mathbf{M}}^{-1}\mathbf{J}^\top$ of rank at most $6N_b$,
and the Woodbury identity turns its inverse into a single
$6N_b \times 6N_b$ Cholesky solve ($24 \times 24$ for Helhest) plus a
row-by-row division by $c_i$. This replaces PCR on up to $170$ rows with an
exact direct solve that runs several times faster in our measurements,
while the solver still takes $\mathbf{A}$ and $\mathbf{b}$ and returns
$\Delta\boldsymbol{\lambda}$ as before.

The engine is built around the constraint-space system
$\mathbf{A}\mathbf{x} = \mathbf{b}$. The Woodbury matrix identity
\citep{woodbury1950inverting,golub2013matrix} gives the body-space solve
directly for this system, without going back to \eqref{eq:kkt}. For
invertible $\mathbf{B}$ and $\mathbf{D}$, it states
\begin{equation}
  \left(\mathbf{B} + \mathbf{U}\mathbf{D}\mathbf{V}\right)^{-1}
  = \mathbf{B}^{-1} - \mathbf{B}^{-1}\mathbf{U}
    \left(\mathbf{D}^{-1} + \mathbf{V}\mathbf{B}^{-1}\mathbf{U}\right)^{-1}
    \mathbf{V}\mathbf{B}^{-1} ,
\end{equation}
whenever the inner matrix is invertible. It inverts an easy matrix plus a
low-rank term, and needs only an inverse of the size of the low-rank term.
With $\mathbf{B} = \mathbf{C}$, $\mathbf{U} = \mathbf{J}$,
$\mathbf{D} = \tilde{\mathbf{M}}^{-1}$ and $\mathbf{V} = \mathbf{J}^\top$,
the left-hand side is $\mathbf{A}^{-1}$:
\begin{equation}
  \mathbf{A}^{-1}
  = \left(\mathbf{C} + \mathbf{J}\tilde{\mathbf{M}}^{-1}\mathbf{J}^\top\right)^{-1}
  = \mathbf{C}^{-1} - \mathbf{C}^{-1}\mathbf{J}\,\mathbf{K}^{-1}\mathbf{J}^\top\mathbf{C}^{-1} .
  \label{eq:woodbury-A}
\end{equation}
$\mathbf{A}$ is a diagonal matrix plus a term of rank at most $6N_b$, and
the only inverse left on the right is the body-space matrix $\mathbf{K}$ of
\eqref{eq:body-system}. Applied to $\mathbf{b}$, \eqref{eq:woodbury-A}
becomes two steps:
\begin{equation}
  \mathbf{K}\mathbf{v} = \mathbf{J}^\top\mathbf{C}^{-1}\mathbf{b},
  \qquad
  \mathbf{x} = \mathbf{C}^{-1}\left(\mathbf{b} - \mathbf{J}\mathbf{v}\right) .
  \label{eq:woodbury-solve}
\end{equation}
The first step is one Cholesky solve with $\mathbf{K}$. The second is a
row-by-row division.

This is the body-space reduction of Section~\ref{sec:body-space-idea},
written for the right-hand side $\mathbf{b}$ of \eqref{eq:ostrich-system}.
The auxiliary vector is $\mathbf{v} = \tilde{\mathbf{M}}^{-1}\mathbf{J}^\top\mathbf{x}$,
the velocity change caused by the constraint forces. Indeed,
$\mathbf{K}\tilde{\mathbf{M}}^{-1}\mathbf{J}^\top\mathbf{x}
= \mathbf{J}^\top\mathbf{C}^{-1}(\mathbf{C} + \mathbf{J}\tilde{\mathbf{M}}^{-1}\mathbf{J}^\top)\mathbf{x}
= \mathbf{J}^\top\mathbf{C}^{-1}\mathbf{b}$.
Hence $\mathbf{x} = \Delta\boldsymbol{\lambda}$ and
$\Delta\mathbf{u} = \mathbf{v} - \tilde{\mathbf{M}}^{-1}\mathbf{r}_u$.
The eliminating-$\Delta\boldsymbol{\lambda}$ view is the physics form; the
Woodbury identity is the linear-algebra form of the same computation. The
Woodbury form keeps the interface of the existing solver: it takes
$\mathbf{A}$ and $\mathbf{b}$ as given and returns $\mathbf{x}$, so it can
replace PCR without changing the rest of the Newton iteration. Because the
adjoint solve uses the same $\mathbf{A}$ (Section~\ref{sec:solvers}), it can
reuse the same Cholesky factor of $\mathbf{K}$.

\subsection{Better Gradients for Differentiable Simulation}
\label{sec:body-space-adjoint}

\paragraph{Differentiable simulation.}
A differentiable simulator returns not only a trajectory but also the
gradient of a scalar loss $L$ on that trajectory with respect to the
initial state, the controls or the model parameters. Trajectory
optimization, system identification and policy learning all use this
gradient. In ostrich a time step is defined implicitly, by a residual
$\mathbf{R}(\mathbf{s}^+, \mathbf{s}^-) = \mathbf{0}$ that Newton's method
solves for the new state $\mathbf{s}^+$. The gradient through one step
follows from the implicit function theorem. Computing it for a scalar loss
takes one linear solve with the transposed Newton matrix, the
\emph{adjoint solve}, whose solution is the \emph{adjoint vector}. Going
backwards through the time steps, each step turns the gradient with respect
to its own state into the gradient with respect to the previous state. Like
the forward step, the adjoint solve reduces to a system with the same
matrix $\mathbf{A}$ \citep[eq.~(18)]{kucera2026ostrich}.

\paragraph{The adjoint step in ostrich today.}
Let $\mathbf{f}$ be the gradient of the loss with respect to the new body
velocities (including the contribution through the poses). The engine
computes the gradient $\mathbf{w}_u$ handed to the previous step in three
steps:
\begin{equation}
  \mathbf{d} = \tilde{\mathbf{M}}^{-1}\mathbf{f}, \qquad
  \mathbf{A}\mathbf{w} = \mathbf{J}\mathbf{d} \;\;\text{(PCR)}, \qquad
  \mathbf{w}_u = \mathbf{d} - \tilde{\mathbf{M}}^{-1}\mathbf{J}^\top\mathbf{w} .
  \label{eq:adjoint-pcr}
\end{equation}
Here $\mathbf{d}$ is the response of the bodies without constraints, and
$\mathbf{w}$ is the constraint part of the adjoint, solved by PCR with the
same tolerance and iteration cap as the forward solve. The last step is
fragile. When a body is held firmly by stiff joints or sticking contacts,
the constraints absorb almost all of $\mathbf{f}$. Then $\mathbf{w}_u$ is
small, but it is computed as the difference of two large, nearly equal
vectors, so a relative error in $\mathbf{w}$ is amplified (catastrophic
cancellation). For one body of mass $m$ and one constraint row with
$J = 1$ and compliance $c$, a relative error $\varepsilon$ in $w$ gives
the relative error $\varepsilon/(mc)$ in $w_u$. The amplification is the
ratio of the constraint stiffness $1/c$ to the inertia $m$. A more accurate
solve for $\mathbf{w}$ reduces $\varepsilon$, but the amplification stays.

\paragraph{The adjoint step in body space.}
Writing \eqref{eq:adjoint-pcr} as one formula gives
\begin{equation*}
  \mathbf{w}_u = \left(\tilde{\mathbf{M}}^{-1} - \tilde{\mathbf{M}}^{-1}\mathbf{J}^\top\mathbf{A}^{-1}\mathbf{J}\tilde{\mathbf{M}}^{-1}\right)\mathbf{f} .
\end{equation*}
The bracket is the right-hand side of the Woodbury identity with
$\mathbf{B} = \tilde{\mathbf{M}}$, $\mathbf{U} = \mathbf{J}^\top$,
$\mathbf{D} = \mathbf{C}^{-1}$ and $\mathbf{V} = \mathbf{J}$, i.e.\ the
identity of Section~\ref{sec:woodbury} used in the other direction, with
$\tilde{\mathbf{M}}$ as the easy matrix. Hence
\begin{equation}
  \mathbf{K}\mathbf{w}_u = \mathbf{f}, \qquad
  \mathbf{w} = \mathbf{C}^{-1}\mathbf{J}\mathbf{w}_u ,
  \label{eq:adjoint-body}
\end{equation}
where the second formula follows from
$\mathbf{C}\mathbf{w} = \mathbf{J}\mathbf{d} - \mathbf{J}\tilde{\mathbf{M}}^{-1}\mathbf{J}^\top\mathbf{w} = \mathbf{J}\mathbf{w}_u$.
Nothing is subtracted: a firmly held body contributes a large stiffness to
$\mathbf{K}$, so $\mathbf{K}^{-1}\mathbf{f}$ comes out small directly. The
adjoint step also costs almost nothing. The forward solve has already
computed the Cholesky factor of $\mathbf{K}$, so \eqref{eq:adjoint-body}
needs only two triangular solves, $\mathcal{O}((6N_b)^2)$ operations.

\paragraph{Comparison.}
We measured the error of a single adjoint step offline. For every
recorded time step we take the matrix of its last Newton iteration
($25$ steps per scene) and $8$ random loss gradients $\mathbf{f}$ on all
bodies. All solvers run in single precision and are compared against a
dense solve in double precision. Table~\ref{tab:adjoint-accuracy} shows
the relative error of $\mathbf{w}_u$.

% Data: step 2d of the deflation study (docs/deflation-study/deflation_ab.md,
% "Adjoint accuracy"; test_scripts/deflation/adjoint_accuracy.py), measured
% 2026-10-13 on the dataset of that date, before the regeneration used in
% Section 4. Random loss gradients only; the chassis-only loss gives the same
% picture.
\begin{table}[h]
  \centering
  \small
  \begin{tabular}{@{}lcccccc@{}}
    \toprule
    & \multicolumn{3}{c}{obstacle scene} & \multicolumn{3}{c}{terrain scene} \\
    \cmidrule(lr){2-4} \cmidrule(l){5-7}
    Adjoint solver & p50 & p90 & $>10\%$ & p50 & p90 & $>10\%$ \\
    \midrule
    PCR Jacobi, tol $10^{-5}$, cap $16$ & $5 \cdot 10^{-2}$ & $0.57$ & $42\%$ & $4 \cdot 10^{-2}$ & $0.16$ & $22\%$ \\
    PCR Jacobi, tol $10^{-3}$, cap $26$ & $4 \cdot 10^{-2}$ & $0.46$ & $39\%$ & $8 \cdot 10^{-3}$ & $6 \cdot 10^{-2}$ & $3\%$ \\
    PCR joint/control, tol $10^{-5}$, cap $16$ & $6 \cdot 10^{-3}$ & $0.15$ & $14\%$ & $1 \cdot 10^{-4}$ & $8 \cdot 10^{-3}$ & $0.5\%$ \\
    Body-space $\mathbf{w}$, then \eqref{eq:adjoint-pcr} & $0.14$ & $2.1$ & $51\%$ & $1 \cdot 10^{-3}$ & $0.28$ & $18\%$ \\
    Body space, \eqref{eq:adjoint-body} & $\mathbf{2 \cdot 10^{-4}}$ & $\mathbf{2 \cdot 10^{-3}}$ & $\mathbf{0\%}$ & $\mathbf{3 \cdot 10^{-6}}$ & $\mathbf{4 \cdot 10^{-4}}$ & $\mathbf{0\%}$ \\
    \bottomrule
  \end{tabular}
  \caption{Relative error of the gradient $\mathbf{w}_u$ after one adjoint
    step, single precision, $25$ time steps $\times$ $8$ random loss
    gradients per scene: median (p50), 90th percentile (p90) and share of
    draws with more than $10\%$ error. The first two rows are the ostrich
    solver at the default and the scene configuration.}
  \label{tab:adjoint-accuracy}
\end{table}

With PCR, as ostrich uses it today, the gradient of one step is often
substantially wrong. In the obstacle scene about $40\%$ of the draws have
more than $10\%$ error, and the worst cosine between the computed and the
exact $\mathbf{w}_u$ is $0.61$ to $0.72$. The terrain scene is milder.
Joint/control deflation (Section~\ref{sec:z-joint-control}) helps, but
does not remove the large errors. The fourth row isolates the cause. It
solves for $\mathbf{w}$ accurately in body space (relative error about
$3 \cdot 10^{-3}$) but keeps the subtraction of \eqref{eq:adjoint-pcr}, and
it is the worst row: in the worst case the gradient even points the wrong
way (cosine $-0.51$). The formula, not the solver, limits the accuracy. The
body-space adjoint \eqref{eq:adjoint-body} has at most $0.5\%$ error on
every draw and a worst cosine of $1.0000$. Together with
Section~\ref{sec:body-space-idea}, the same Cholesky factor of $\mathbf{K}$
thus gives a faster forward solve and an essentially exact adjoint step at
almost no extra cost.

These are errors of a single step, measured offline on the matrix of the
last Newton iteration, which approximates the converged state. Over a
trajectory the errors propagate through all earlier steps, which probably
increases the difference, but this has not been measured. The decisive
tests, a finite-difference check of the gradient and a trajectory
optimization with both adjoints, require implementing the body-space solve
in the engine.
